% What this holds: Section Board Structure (Duties, Seats).
% Read by arlington-bsap.tex. Prose lives here; formatting lives in the wrapper and style/.
\section{Board Structure}

\subsection{Board Duties}

From 1847 a county court of elected justices of the peace governed Alexandria
County, judging its cases and, once a year, levying its
taxes.\autocite[ch.~53, secs.~1 and 3]{vacode1860} Until 1870 the county had
no elected body whose only business was to govern it. The board of supervisors
was that body, one member from each township: the constitution had it meet
each December to audit the county's accounts, equalize the valuation of
property and fix the county levy for the year, and left the General Assembly to
add to its duties by statute.\autocite[art.~VII, sec.~2]{vaconstitution1869}

Three supervisors were the county's elected government until November 4, 1930,
when voters adopted the County Manager Plan, 1,936 to 428.\autocite[194--196]{rose1976}
The Plan put a second office beside the Board. The Board appoints a County
Manager, who holds the county's administrative and executive powers, including
the appointment of officers and employees, and the Board may not move any budget
allocation by more than 10 percent.\autocite{vaacts1930c167}

The act left one question open, and it was the one that mattered: who appoints
the department heads. The Board had made those appointments until 1937, when it
delegated the power to the Manager. In January 1951 it took the power back by
resolution, and in March 1951 a circuit judge ruled that it had no right to give
the power up.\autocite{dailysun1951spicer,sun1951rulingboard} The next spring the
General Assembly barred a Board member from directing or requesting an
appointment below the Board, on pain of losing his
seat,\autocite{vaacts1952c443} and in November 1952 voters settled the matter,
the last time they amended the plan, by approving both questions the General
Assembly put to them: that the Manager serve an indefinite term, removable at
any time by the Board, and that the Manager appoint department
heads.\autocite{vaacts1952c198,dailysun1952results} From then on the Board hired
and could dismiss one executive and could not reach past him to the staff.

The County's own rules have since applied that division to particular jobs.
Under the County Code the Manager appoints and removes employees in the
competitive service and the executive management service and writes the
personnel regulations, while the Board approves the pay and classification plans
and may not set the pay of any one competitive
employee.\autocite{arlingtoncode6civilservice} The Board appropriates money by
department, and the Manager may move funds within a
department.\autocite{arlingtonva2026budgetresolution} A capital construction
contract above \$1,000,000 needs the Board's
approval,\autocite{arlingtonva2025purchasingresolution} and even the Board's
agenda passes through the Manager, who drafts the list of proposed items about
two weeks before each regular meeting for the Chair to
approve.\autocite{arlingtonva2026procedures} The Board, in short, sets the
budget, the pay plan and the large contracts, and the Manager runs the staff
beneath them.

The Board keeps a few officers of its own. A 1968 statute let any county create
the office of county attorney, serving at the governing body's
pleasure.\autocite{vaacts1968c695} The General Assembly has since let a board under
the plan appoint its own county auditor, in 2015, and an independent policing
auditor, in 2026, each serving at the board's
pleasure.\autocite{vacode1527092,vaacts2026c372} Arlington's Board hires its
policing auditor under an ordinance effective July 1,
2026.\autocite{arlingtoncode69oversight} Where the state's law, not
the County's, limits what the Board may do, as in \textit{County Board of
Arlington v. Brown} (1985), Section~\ref{sec:legal} takes it up.

\subsection{Board Seats}

For its first 60 years Arlington was a small county that grew slowly, from about
3,000 residents in 1870 to 27,000 in 1930. Then federal employment arrived, first
under the New Deal and then with the militarization of the 1940s, and the county
reached 57,000 residents in 1940 and 135,000 in 1950, five times its size 20
years earlier.\autocite[239--40, 248]{bestebreurtje2017} It has not stopped: the
2020 census counted 239,000.

\begin{pairpage}
  \refstepcounter{figure}\label{fig:residents-per-seat}
  \figuretitle{County Population and Board Size}
  \medskip
  \pairfigure{residents_per_seat.pdf}
  \smallskip
  \figurenotes{
    }{Source: US decennial census, 1870--2020.}
  \vfill
  \refstepcounter{figure}\label{fig:residents-district-map}
  \figuretitle{The Three Magisterial Districts, 1870--1931}
  \smallskip
  \includegraphics{residents_by_district_map.pdf}
  \vspace{2pt}
  \figurenotes{
    Notes: The county as it stood before Alexandria annexed land from it in 1915
    and 1930, with the three districts as they stood from 1870 until the Board
    went at-large in 1932. The 1907 county map is the only map found that draws
    the district lines, and the lines here are read from it to within a few
    hundred meters. The paler areas are the land Alexandria annexed in 1915 and
    1930, almost all of it Jefferson district. The city as it stood in 1912 is
    left blank. The Potomac is the right-hand edge.}{Source: Noetzel and Boteler,
    Map of Alexandria County (1907); Rose (1964); Census Bureau and Arlington
    County boundary files.}
\end{pairpage}

The Board has changed size once. For those first 60 years it was three
supervisors, one each from the Arlington, Jefferson and Washington districts
(Figure~\ref{fig:residents-district-map}), and
the referendum of November 4, 1930, that adopted the Plan also replaced them with
five members elected at large (Section~\ref{sec:method}); the first such Board
took its seats in January 1932. With the population of 1930, the change cut the
residents each seat stood for from about 9,000 to about 5,000. The county then
outgrew it within a decade, to 11,000 residents a seat in 1940 and 27,000 in
1950, and each seat stands for about 48,000 today, against about 1,000 in 1870.

The General Assembly did offer a way to grow the Board. In 1952 it let Arlington
voters put a seven-member Board on the ballot by a petition of 200 voters filed by
August 1, and nobody filed one.\autocite{vaacts1952c591} It reopened the window in
1954 and 1958, for the question of election every two years only, and again
nobody petitioned.\autocite{vaacts1954c151,vaacts1958c207} Part C sets
Arlington's residents per member beside those of Virginia's other large
localities (Figure~\ref{fig:localities-peers}).

